\section{L1 Foundations}
\subsection*{Token / Type / Morphology}
Type: unique vocabulary item in $V$; token: one occurrence; TTR $=|V|/N$.\\
Heaps/Herdan law: $|V|\approx KN^{\beta},\ \beta\approx 0.5$ (type growth slows as corpus grows).\\
Morpheme: minimal meaning unit; root + affix; inflection vs derivation.\\
Orthographic token and grammatical word may differ (e.g., \texttt{I'm}).

\subsection*{Unicode and UTF-8}
Unicode maps symbols to code points (e.g., \texttt{a}=U+0061); code point $\neq$ glyph.\\
UTF-8 patterns: 1-byte $0xxxxxxx$, 2-byte $110xxxxx\ 10xxxxxx$, 3-byte $1110xxxx\ 10xxxxxx\ 10xxxxxx$, 4-byte $11110xxx\ 10xxxxxx\ 10xxxxxx\ 10xxxxxx$.\\
ASCII stays 1 byte in UTF-8; Python 3 \texttt{len()} counts code points, not bytes.

\subsection*{BPE Essentials}
Train ($k$ merges): frequent pair $(A,B)\to AB$, add to $V$, replace all $A\ B$.\\
Encode: start from chars/bytes, then apply learned merges in learned order (greedy).\\
Byte-level BPE uses 256-byte alphabet and avoids byte-level OOV.

\subsection*{Regex Quick Sheet}
Classes \verb|[A-Z]|, \verb|[^0-9]|; quantifiers \verb|* + ? {m,n}|; anchors \verb|^ $ \b|.\\
Groups \verb|(...)| / \verb|(?:...)|; lookahead \verb|(?=...)| / \verb|(?!...)|; substitution via \texttt{re.sub()}.

\subsection*{L1 Reminders}
Tokenizer choice (especially CJK segmentation) strongly changes token budget.\\
Perplexity comparison is valid only under same tokenizer, preprocessing, and test split.
